In this section, the result that monomial ideals have finite generators is extended to polynomial ideals. The basic idea is that by rearranging the ideal to a monomial ideal using the leading term, we can apply the Dicksons lemma. The first step towards bridging the gap between polynomial and monomial ideals is by crafting an ideal of interest. Namely, the ideal generated by the leading term of a polynomial ideal. 

\begin{definition}
    Let $I \subseteq \polyRing$ be an ideal other than $\lbrace 0\rbrace$ and fix a monomial ordering. Then 
    \begin{enumerate}
        \item $LT(I) = \{cx^\alpha$ | there exists $f \in I \backslash \{0\}$ with $LT(f) = cx^\alpha\}$.
        \item $\ideal{LT(I)}$ is the ideal generated by the elements of $LT(I)$.
    \end{enumerate}
\end{definition}

$LT(I)$ is a large set, nut it is the bridge between general ideals and monomial ideals. After all, the difference between a monomial and the lead term of the polynomial are the other monomials in the polynomial. If we can capture them by constructing an ideal that generates everything by using the leading term, we can create a superset to operate upon. $\ideal{LT(I)}$ is that set.  

\begin{proposition}
Let $I \subseteq \polyRing$ be an ideal different from the  $\lbrace 0\rbrace$.
\begin{enumerate}
    \item $\ideal{LT(I)}$ is a monomial ideal.
    \item There are $\vectorComponentsX{g}{t} \in I$ such that
          $\ideal{LT(I)} = \ideal{LT(g_1), \dots, LT(g_t)}$.
\end{enumerate}
\end{proposition}

\begin{proof}
    Since we have used the leading term to construct an ideal, the difference between the monomial ideal and $\ideal{LT(I)}$ is a nonzero constant. Consider a single $g_i = g$, and break the lead term down to the lead coefficient and lead monomial. $LT(g) = LC(g) \cdot LM(g)$ means that $LT(g)$ is a field scalar of $LM(g)$, which means that $LT(g)\in \ideal{LM(g)}$. We can multiply by the inverse to show that $LT(g) \cdot LC^{-1}(g) = LM(g)$. But $LC^{-1}(g)$ is still in $\polyRing$, so the same logic applied and they are equal to each other. So, $\ideal{LT(I)} = \ideal{LM(I)}$.
    The same logic is used for the second item. Using Dicksons Lemma tells us that monomial ideals are finitely generated. Since $\ideal{LT(I)}$ only differs from a monomial ideal by a nonzero constant, Dicksons lemma can be applied here as well and $I$ is finitely generated.
\end{proof}

With all this information accumulated, we can formally prove that every polynomial ideal has a finite generating set. This theorem is known as the Hilber Basis Theorem.

\begin{theorem}[Hilber Bases Theorem]
   Every ideal $I \in\polyRing $ has a finite generating set. In other words, $ I = \ideal{\vectorComponentsX{g}{t}}$ for some $\vectorComponentsX{g}{t} \in I$.
\end{theorem}

This proof is outside of the scope, but is given in chapter 2 of Cox \cite{Theory} as Theorem 4 on page 77. 